\clearpage\corrigesChapitre{10}{Transformations nucléaires}
\corrigesSouples
\begin{blocLivre}{cor-10-1}\begin{corrige}{1}{Reconnaître des isotopes}
\begin{questions}
\item Les deux noyaux d’oxygène ont le même $Z=8$ et des $A$ différents.
\item Ils possèdent 8 et 10 neutrons. Le fluor en possède 10, mais son $Z$ est différent.
\end{questions}
\end{corrige}\end{blocLivre}
\begin{blocLivre}{cor-10-2}\begin{corrige}{2}{Classer des transformations}
\begin{questions}
\item Vaporisation : physique ; combustion : chimique ; fission : nucléaire.
\item La fission.
\end{questions}
\end{corrige}\end{blocLivre}
\begin{blocLivre}{cor-10-3}\begin{corrige}{3}{Un neutron}
\begin{questions}
\item Il compte pour un nucléon et pour zéro proton.
\item Sa charge électrique est nulle.
\end{questions}
\end{corrige}\end{blocLivre}
\begin{blocLivre}{cor-10-4}\begin{corrige}{4}{La fusion de la glace}
\begin{questions}
\item C’est une transformation physique, un changement d’état solide-liquide.
\item Non, elles restent des molécules de \ce{H2O}.
\end{questions}
\end{corrige}\end{blocLivre}
\begin{blocLivre}{cor-10-5}\begin{corrige}{5}{Une fusion fournie}
\schema[corrige]{nucleaire}
\begin{questions}
\item $A : 2+3=4+1=5$ ; $Z : 1+1=2+0=2$.
\item Deux noyaux légers donnent un noyau plus lourd, avec émission d’un neutron.
\end{questions}
\end{corrige}\end{blocLivre}
\begin{blocLivre}{cor-10-6}\begin{corrige}{6}{Compléter une fission}
\begin{questions}
\item $235+1=141+92+x$, donc $x=3$.
\item $92=56+36$ ; les neutrons n’ajoutent pas de charge.
\end{questions}
\end{corrige}\end{blocLivre}
\begin{blocLivre}{cor-10-7}\begin{corrige}{7}{La source d’énergie du Soleil}
\begin{questions}
\item À des transformations nucléaires de fusion dont le bilan convertit de l’hydrogène en hélium.
\item Une combustion réorganise les atomes sans transformer les noyaux ; ce n’est pas le mécanisme énergétique principal du Soleil.
\end{questions}
\end{corrige}\end{blocLivre}
\begin{blocLivre}{cor-10-8}\begin{corrige}{8}{Une chaîne de conversions}
\begin{questions}
\item Les réactions chauffent un fluide ; un dispositif met en mouvement une turbine entraînant un alternateur.
\item Une part est transférée à l’environnement ; les conversions ne dirigent pas toute l’énergie vers la sortie électrique utile.
\end{questions}
\end{corrige}\end{blocLivre}
\begin{blocLivre}{cor-10-9}\begin{corrige}{9}{Un noyau inconnu}
\begin{questions}
\item $A_X=14+4-17=1$ ; $Z_X=7+2-8=1$.
\item X est un proton, noyau $\ce{^{1}_{1}H}$.
\end{questions}
\end{corrige}\end{blocLivre}
\begin{blocLivre}{cor-10-10}\begin{corrige}{10}{Un bilan énergétique fourni}
\begin{questions}
\item $E=3,2\times10^{-11}\times2,0\times10^{18}=\qty{6,4e7}{J}$.
\item Chaque événement est supposé libérer la même énergie dans le bilan choisi.
\end{questions}
\end{corrige}\end{blocLivre}
\begin{blocLivre}{cor-10-11}\begin{corrige}{11}{Charge et nucléons}
\begin{questions}
\item La charge : $10\neq4+5$.
\item Il faut respecter simultanément les lois de conservation adaptées à toutes les particules de l’équation.
\end{questions}
\end{corrige}\end{blocLivre}
\begin{blocLivre}{cor-10-12}\begin{corrige}{12}{Isotopes et ions}
\demarche{Comparer les noyaux avant d’examiner le cortège électronique.}
\begin{questions}
\item Non : leur noyau est le même, avec les mêmes $A$ et $Z$.
\item L’ion a un électron de moins. Une ionisation concerne le cortège électronique.
\end{questions}
\retiens{Une ionisation ne crée pas un nouvel isotope.}
\end{corrige}\end{blocLivre}

